\documentclass{article}

\usepackage[utf8]{inputenc}
\usepackage[portuguese]{babel}
\usepackage[T1]{fontenc}
\usepackage{lmodern}
\usepackage{fullpage}
\usepackage[usenames]{color}
\usepackage[color]{coqdoc}
\usepackage{url}
\usepackage{makeidx,hyperref}
\usepackage[all]{xy}

\usepackage{mathpartir}
\usepackage{tcolorbox}
\usepackage{amsmath,amssymb}

\newcommand{\tto}{\twoheadrightarrow}
\newcommand{\ott}{\twoheadleftarrow}

\title{Formalização da correção do algoritmo Bubble Sort}
\author{%
  Danilo Silveira --- 222014142 \\
  Sofia Mascarenhas Ikawa --- 211060700
}
\date{\today}

\begin{document}
\maketitle

\section*{Participantes}

\begin{itemize}
  \item Danilo Silveira (222014142): lemas \texttt{bubble\_perm},
        \texttt{bs\_permuta} e teorema \texttt{bs\_correto}.
  \item Sofia Mascarenhas Ikawa (211060700): lemas \texttt{bubble\_sorted},
        \texttt{bubble\_min}, \texttt{bubble\_insert\_sorted} e \texttt{bs\_sorted}.
\end{itemize}

\section*{Lema \texttt{bubble\_perm}}

Este lema estabelece que, para toda lista \texttt{l}, aplicar \texttt{bubble l}
produz uma permutação de \texttt{l}. Ou seja, uma passada de borbulhamento
reordena os elementos, mas não cria nem remove nenhum deles.

A prova utiliza indução funcional sobre \texttt{bubble l}, seguindo a mesma
estratégia empregada em \texttt{bubble\_sorted}.
Como \texttt{bubble} foi definida com \texttt{Function}, a indução gera quatro
casos, correspondentes aos ramos da função:

\begin{itemize}
  \item \textbf{Casos base} (\texttt{nil} e \texttt{x::nil}): a saída coincide
        com a entrada; a prova termina com \texttt{reflexivity}.
  \item \textbf{Caso \texttt{x <= y}}: \texttt{bubble} retorna
        \texttt{x :: bubble (y :: l0)}. Aplicamos \texttt{perm\_skip} e a
        hipótese indutiva \texttt{IHl0}, que garante
        \texttt{Permutation (y :: l0) (bubble (y :: l0))}.
  \item \textbf{Caso \texttt{x > y}}: \texttt{bubble} retorna
        \texttt{y :: bubble (x :: l0)}. Como o primeiro elemento muda, usamos
        \texttt{perm\_trans} com a lista intermediária \texttt{y :: x :: l0}:
        \texttt{perm\_swap} (com \texttt{Permutation\_sym}) formaliza a troca
        local de \texttt{x} e \texttt{y}; em seguida, \texttt{perm\_skip} com
        \texttt{IHl0} conclui o caso.
\end{itemize}

\section*{Lema \texttt{bs\_permuta}}

Este lema mostra que o algoritmo completo \texttt{bs} preserva permutação:
\texttt{forall l, Permutation l (bs l)}.

A prova é por indução estrutural sobre a lista.
No \textbf{caso base}, \texttt{bs nil = nil} e a conclusão é imediata
(\texttt{simpl} e \texttt{reflexivity}).

No \textbf{caso indutivo} \texttt{a :: l}, pela definição de \texttt{bs} temos
\texttt{bs (a :: l) = bubble (a :: bs l)}.
O objetivo é provar \texttt{Permutation (a :: l) (bubble (a :: bs l))}.
Utilizamos \texttt{perm\_trans} com a lista intermediária \texttt{a :: bs l}:

\begin{enumerate}
  \item \texttt{Permutation (a :: l) (a :: bs l)} --- obtido pela hipótese
        indutiva (\texttt{Permutation l (bs l)}) estendida com \texttt{perm\_skip}.
  \item \texttt{Permutation (a :: bs l) (bubble (a :: bs l))} --- obtido
        diretamente pelo lema \texttt{bubble\_perm}.
\end{enumerate}

A transitividade encadeia os dois passos e fecha o caso indutivo.

\section*{Teorema \texttt{bs\_correto}}

O teorema \texttt{bs\_correto} é o resultado principal do enunciado:
\texttt{forall l, Sorted le (bs l) /\ Permutation l (bs l)}.
Ele afirma que \texttt{bs} retorna uma lista ordenada que é permutação da
entrada, ou seja, que o algoritmo está correto.

A prova separa as duas propriedades com o tático \texttt{split}:
\begin{itemize}
  \item \texttt{Sorted le (bs l)} --- concluído aplicando o lema \texttt{bs\_sorted}
        (Sofia).
  \item \texttt{Permutation l (bs l)} --- concluído aplicando o lema
        \texttt{bs\_permuta} (Danilo).
\end{itemize}

Nenhuma indução adicional é necessária, pois ambos os lemas já haviam sido
estabelecidos anteriormente.

\section*{Conclusão}

Formalizamos a correção do bubble sort no Coq, demonstrando que a função
\texttt{bs} produz uma lista ordenada e permutação da entrada.
As provas de ordenação foram desenvolvidas por Sofia, enquanto as provas de
permutação e o teorema final foram desenvolvidos por Danilo.
Todos os resultados exigidos pelo enunciado terminam em \texttt{Qed}, e o
projeto compila com \texttt{make coq}.

\input{bubble_sort.v}

\end{document}
